% Fragmento público generado desde propietarios íntegros.
% Residencia: 52.3.1.
% Fuente: ../incoming/radion_monografia_20260827/manuscrito/sections/04_cincuenta_critica.tex:1-105
\noindent La coordenada crítica enlaza la carta de Cayley, el registro APP y la fase dodecafásica mediante aplicaciones tipadas.\par\medskip

\subsubsection{Dos realizaciones exactas de la coordenada angular \(50^\circ\)}

\paragraph{Segunda fase del sistema temporal dodecafásico}

Por \eqref{eq:rueda-dodecafasica}, la sucesión de centros comienza
\[
20^\circ,\ 50^\circ,\ 80^\circ,\ 110^\circ,\ldots.
\]
Resolver
\[
30m-10=50
\]
da únicamente \(m=2\). Por tanto, \(50^\circ\) no es una redondeo de \(A\),
una media vuelta ni el centro de un círculo: es una fase tipada del calendario
dodecafásico. La fracción de vuelta es
\[
\frac{50}{360}=\frac5{36},
\]
exactamente el primer sumando de \(S_E/T_+\).

\paragraph{Igualación de la coordenada crítica de Cayley}

La involución funcional
\[
\rho\longmapsto1-\bar\rho
\]
fija
\[
\operatorname{Fix}(\rho\mapsto1-\bar\rho)
=\{\rho:\Re\rho=1/2\}.
\]
El número \(1/2\) es la coordenada real autodual del intervalo funcional
\(0\leq\Re s\leq1\). No es \enquote{la mitad de una recta} en sentido
longitudinal, y no contiene la altura espectral \(\gamma\). Para
\[
\rho_\gamma=\frac12+i\gamma,
\]
la transformación de Cayley se puede escribir
\begin{equation}
\tau_\gamma
=\frac{\rho_\gamma-1}{\rho_\gamma}
=\frac{\gamma+i/2}{\gamma-i/2}
\in S^1,
\qquad
\gamma=\frac12\cot\frac{\theta_\gamma}{2}.
\label{eq:cayley-critica}
\end{equation}
Así, \(1/2\) fija la costura, \(\gamma\) determina el carácter angular y
\(\Gamma_9\) conserva la profundidad de sus iteraciones.

\paragraph{Inyección afín \(j_{100}\) del registro APP en la carta angular}

La completación cúbica dispone una cámara de diez marcas por coordenada. Una
cara contiene \(10^2=100\) posiciones. Declaramos la carta afín
\begin{equation}
j_{100}:[0,1]\longrightarrow[0,100^\circ],
\qquad j_{100}(\sigma)=100\sigma^\circ.
\label{eq:j100}
\end{equation}
Entonces
\[
j_{100}\left(\frac12\right)=50^\circ.
\]
Al exigir simultáneamente \(j_{100}(1/2)=\theta_m\), resulta
\[
100\cdot\frac12=30m-10
\quad\Longleftrightarrow\quad m=2.
\]

\begin{theorem}[Igualador crítico--dodecafásico]
En la carta APP de cien marcas, la coordenada fija de la involución crítica
se publica exactamente en la segunda fase de la rueda:
\begin{equation}
\boxed{j_{100}(1/2)=\theta_2=50^\circ.}
\label{eq:igualador-50}
\end{equation}
La fase \(m=2\) es única.
\end{theorem}

El teorema es exacto una vez declarada \eqref{eq:j100}. Su estatuto es
\status{formalización nueva}: el medio crítico, la completación \(10^3\), la
cara de cien marcas y \(\theta_2\) ya estaban construidos; la elección de esa
cara como carta crítica afín reúne por primera vez sus dominios.

\paragraph{Compatibilidad del valor crítico con la carta de Cayley}

La transformación de Cayley no envía \(1/2\) a \(50^\circ\). Si
\(\gamma=0\), \eqref{eq:cayley-critica} da
\[
\tau_0=-1=e^{i\pi},
\]
es decir, \(180^\circ\). Los \(50^\circ\) pertenecen a la carta APP de cien
marcas y al calendario dodecafásico; el ángulo Cayley depende de \(\gamma\).
Esta distinción impide colapsar todos los ceros de la función zeta en una sola
fase.

\begin{falsifier}
La identificación \(1/2\leftrightarrow50^\circ\) falla si se elimina la
carta \eqref{eq:j100}. La igualdad escalar \(100(1/2)=50\) no basta para
identificar por sí sola dos objetos. El teorema depende de la cara APP
distinguida y de la rueda dodecafásica.
\end{falsifier}

